% The compile-ready copy of this manuscript lives in overleaf/main_claude.tex,
% alongside sn-jnl.cls, the bibliographies and the figures. Compile from there;
% the \includesvg paths in this file are relative to the overleaf directory.

\documentclass[pdflatex,sn-mathphys-num]{sn-jnl}

\usepackage{graphicx}%
\usepackage{multirow}%
\usepackage{amsmath,amssymb,amsfonts}%
\usepackage{amsthm}%
\usepackage{mathrsfs}%
\usepackage[title]{appendix}%
\usepackage{xcolor}%
\usepackage{textcomp}%
\usepackage{manyfoot}%
\usepackage{booktabs}%
\usepackage{algorithm}%
\usepackage{algorithmicx}%
\usepackage{algpseudocode}%
\usepackage{listings}%

\usepackage{caption}
\usepackage{svg}
\usepackage{hyperref}
\newcommand{\figpanel}[2]{\hyperref[#1]{Fig.~\ref*{#1}#2}}

\raggedbottom

\begin{document}

%% Header
\title{Estimating Neural Timescales}

\author*[1]{\fnm{Ryan} \sur{Hammonds} \email{rphammonds@ucsd.edu}}
\author[2]{\fnm{Dillan} \sur{Cellier}}
\author[3]{\fnm{Blanca} \sur{Martin-Burgos}}
\author[3]{\fnm{Trevor} \sur{McPherson}}
\author[4]{\fnm{Richard} \sur{Gao}}
\author*[1, 2, 3, 5]{\fnm{Bradley} \sur{Voytek} \email{bradley.voytek@gmail.com}}

\affil[1]{\orgdiv{Halıcıoğlu Data Science Institute}}
\affil[2]{\orgdiv{Department of Cognitive Science}}
\affil[3]{\orgdiv{Neurosciences Graduate Program}}
\affil[4]{\orgdiv{Kavli Institute for Brain and Mind}, \orgname{University of California, San Diego}, \city{La Jolla}, \postcode{92093}, \state{CA}, \country{USA}}

\affil[5]{\orgdiv{Bernstein Center for Computational Neuroscience}}
\affil[6]{\orgdiv{Machine Learning in Science}, \orgname{University of Tübingen \& Tübingen AI Center}, \state{Tübingen}, \country{Germany}}

\abstract{
In the attempt to understand the ``neural code'', we have developed conceptual and mathematical tools for interpreting neuronal activity across recording modalities. Neural activity is predictable from its own recent past, and that predictability fades with time. In sensory cortices, all forms of neural activity, from spiking to EEG, can reliably predict future neural activity up to a few hundred milliseconds later. In contrast, activity in association cortex predicts subsequent activity up to several seconds later. This temporal dependency of neural activity is often referred to as ``neural timescales''. There are several mathematical models for quantifying timescales: extracting a decay parameter from the autocorrelation function, fitting functions to power spectra, or directly from the time series with autoregressive models. Here, we show in both simulated and real neural data how common neural signal features -- model mismatch, neural oscillations, finite sampling, and dynamic timescales -- can yield inaccurate or uninterpretable timescale estimates. We show that higher-order autoregressive, AR($p$), models, designed to predict future data from $p$ past time points, handle these issues. The AR($p$) model decomposes into a set of modes or components. In autocorrelation, these components are represented as a sum of damped oscillators or relaxation functions, both with exponentially decaying structure. This method is applied to a variety of neural recordings from different species, spanning fMRI, MEG, EEG, ECoG, LFP and spiking activity.
}

\maketitle

\section{Introduction}\label{sec:intro}

\subsection{Timescales}\label{sec:timescales}

Many real world processes are correlated at varying scales. For example, when trying to predict tomorrow's high temperature, a reasonable starting point is to use the high temperature from today. This also works to predict the high temperature the day after, and the day after that, and so on. One common approach for capturing these temporal relationships is through the use of an autocorrelation, which quantifies how correlated a signal is with itself at every temporal lag. The autocorrelation function (ACF) describes the resulting curve, which can be visualized as the correlation at each time lag (\figpanel{fig:ar1}{}).

In natural processes the ACF often decays exponentially, with correlation strongest at short time lags and decaying as the lag increases. The rate at which the ACF decays from correlated to uncorrelated is usually described as the ``timescale'' of the process. For purely random, Gaussian distributed, white noise processes, the decay is instant -- dropping from perfectly correlated at lag-zero to uncorrelated at all other lags. Timescales are quantified as the exponential decay rate of the ACF, where numerous real-world processes, ranging from the atomic scale with radioactive decay \cite{semkow_oscillations_2009}, to the astronomical with stellar flares \cite{anfinogentov_decaying_2013}, show this temporal dependence. Regardless of the domain or modality, the mathematical basis of timescales is the same, while interpretations and ground truth mechanisms will vary based on the signal's context.

Returning to the temperature example, temperature fluctuates periodically, on a day/night and seasonal basis. These periodic fluctuations introduce oscillations into the ACF. An oscillation is strongly anticorrelated with itself at every half-cycle lag and strongly correlated at every full-cycle lag. Weather does not generate a perfect sinusoids, so the amplitude and timing of each cycle will vary, and the correlation between cycles weakens as lag increases. The ACF of such a process is a damped oscillator: a cosine at the period of the cycle, multiplied by an exponentially decaying envelope. Two parameters describe this, a period and a decay rate, rather than the single decay rate of a purely aperiodic process. Yule described the sunspot cycle in these terms, as a damped oscillator, and introduced the autoregressive (AR) model to do so \cite{yule_1927}.

In this paper we compare the ACF, PSD, and AR models commonly used to estimate timescales, and show how each is biased when the assumed form is incorrect. This is shown in simulation, where the timescale is known, and in recorded data, where it is not, and recommendations are provided for addressing these biases. Importantly, we show how to interpret complex, high-order AR($p$) models from the perspective of timescales. AR($p$) has a corresponding form in all three domains and can be decomposed into a set of decaying and oscillatory components.

\begin{figure}[htbp]
  \centering
  \includesvg[width=0.95\linewidth]{revision/figures/fig_ar1.svg}
  \caption{Timescales, represented as an AR(1) model and process (left), autocorrelation (middle), and power spectral density (right). A key connection is between the AR(1) coefficient and timescales, where $\varphi$ is equal to the autocorrelation at lag one, $A(1)$, and the knee of the spectrum sits at $1/(2\pi\tau)$. Each column is the basis of one common estimator, and the three estimate the same parameter of the same model.}
  \label{fig:ar1}
\end{figure}

\subsection{Neural Timescales}\label{sec:neural}

In the context of neuroscience, timescales have been observed across recording modalities, from single unit recordings \cite{murray_hierarchy_2014} and neuronal population local field potentials (LFP) \cite{haider_millisecond_2016} to human electroencephalography (EEG) \cite{meisel_decline_2017}, magnetoencephalography (MEG) \cite{zhigalov_relationship_2015}, and functional magnetic resonance imaging (fMRI) \cite{watanabe_atypical_2019, raut_organization_2020}. Neural timescales characterize the length of time in which an input is integrated and maintained by neuronal populations. That is, if an input causes a change in neural activity that then immediately reverts back to random, its timescale will be very short, or fast: you can't predict future values from the current state. However, if an input causes a lasting, predictable structure in activity, its timescale will be longer, or slower: given its current state, you can predict future values. This is interpreted as that activity having a longer ``memory'' of the input. This interpretation of neural timescales as the maintenance of information, or memory, in a neural time series is what makes timescales particularly attractive from a computational neuroscience perspective. However, this is also what makes accurate measurement of timescales so critical, especially given the rapid emergence of studies that examine the functional role of timescales in cognition \cite{cavanagh_diversity_2020, gao_neuronal_2020} and their alteration in disease states \cite{watanabe_atypical_2019, wengler_distinct_2020}.

Related to the potential functional role of neural timescales is the non-random structure of neural timescales across the brain. Rather, they are spatially organized in a cortical hierarchy \cite{hasson_hierarchy_2008, kiebel_hierarchy_2008, murray_hierarchy_2014, gao_neuronal_2020, misic_brain_2010, raut_organization_2020}. The lowest levels of the hierarchy are associated with short timescales, on the order of milliseconds, in visual \cite{buracas_efficient_1998} and auditory \cite{yang_millisecond-scale_2008} sensory cortices. At higher levels along the cortical hierarchy, longer timescales, on the order of seconds, are associated with decision making \cite{brody_timing_2003} and working memory \cite{wang_probabilistic_2002} in the prefrontal and parietal cortices. This hierarchical distribution has also been observed across species, suggesting evolutionary preservation \cite{buzsaki_scaling_2013}. The current work focuses on neural timescales, rather than 2D and 3D spatial extensions, e.g. spatial scales. However, the spatial distribution of timescales means that a timescale should be estimated within a local region that is approximately stationary.

Periodic signals introduce oscillations in the ACF. Neural oscillations are embedded within a non-oscillatory, aperiodic signal, which is thought to reflect asynchronous neuronal drive, capturing the timescale of the underlying neural activity \cite{gao_neuronal_2020}. This mixing of aperiodic and periodic features means that the ACF also blends aperiodic timescales with neural oscillations. As we will explore in more detail below, this mixing of periodic and aperiodic processes can result in large biases and errors in estimates of the timescale from the ACF.

\subsection{Estimation}\label{sec:estimation}

Timescales are estimated as a parameter of a model, and is usually a form of the first-order autoregressive, AR(1), model: $T_{\text{tomorrow}} = \varphi T_{\text{today}} + \epsilon$. Assuming tomorrow's temperature will be very close to today's, the coefficient, $\varphi$, will approach one. As $\varphi$ approaches zero, today's temperature is increasingly less indicative of tomorrow's temperature and trends towards a Gaussian distributed random white noise process dominated by $\epsilon$.

The AR(1) model may be applied to varying temporal scales. If $T_0$ is today's temperature and $T_k$ is the temperature in $k$-days, we can update the temperature model to: $T_k = \varphi^k T_0 + \epsilon$. In these lagged cases, the AR(1) coefficient, $\varphi$, depends on the discrete time step, $k$. This is equivalent to the value of the autocorrelation function (ACF) at the corresponding lag step, given the sampling rate, $f_s$. The ACF, $A(k)$, of an AR(1) process is described by an exponentially decaying function and may be alternatively expressed in terms of the AR(1) coefficient, $\varphi$ (eq. \ref{eq:acf}). The autocorrelation at lag one is equal to $\varphi$. The AR(1) model can simulate a signal, $x$, given a white noise input (eq. \ref{eq:ar}). The ACF and PSD are connected through the Fourier transform, described by the Wiener-Khinchin theorem, so the AR(1) model also has a form in the power spectral density (PSD) (eq. \ref{eq:psd1}).

Timescales, $\tau$, may be expressed as the AR(1) coefficient, $\varphi$ (eq. \ref{eq:tau}), and are bounded by the scale of the observation, defined by the sampling rate, $f_s$, and the length of the signal, $n$. Sampling rate lower bounds observable timescales, as very short timescales, relative to $f_s$, will be undetectable. For example, if a timescale related to temperature fluctuations is on the order of days, but temperature is sampled once a week, the daily timescale will be impossible to detect. Inversely, the length of the signal upper bounds detectable timescales. If a timescale is on the order of years, but only one week of data is collected, the yearly timescale will be undetectable.

\begin{align}
x_t &= \varphi\, x_{t-1} + \epsilon \label{eq:ar} \\
A(k) &= e^{\frac{-k}{\tau f_s}} = \varphi^k \label{eq:acf} \\
\tau &= \frac{-1}{f_s \log{\varphi}} \label{eq:tau} \\
P(f) &= \frac{\sigma^{2}}{\left|1 - \varphi\, e^{-j 2 \pi f / f_s}\right|^{2}} \label{eq:psd1}
\end{align}

Neural processes often contain complex structure, beyond what an AR(1) model can sufficiently describe. A solution is increasing complexity through AR($p$) models that use the last $p$ values to predict the next, with one coefficient per lag (eq. \ref{eq:arp}). The ACF and PSD also have corresponding AR($p$) forms. Every past value the model uses adds a term to the ACF, so that the correlation at each lag is built from the correlations at the $p$ lags before it (eq. \ref{eq:arp_acf}), and adds a term to the PSD (eq. \ref{eq:arp_psd}). This lets both take shapes an AR(1) cannot, such as an ACF that dips below zero before returning, or a spectrum with both periodic and aperiodic components.

\begin{align}
x_t &= \sum_{k=1}^p \varphi_k x_{t-k} + \epsilon \label{eq:arp} \\
A(k) &= \sum_{m=1}^p \varphi_m A(k-m), \qquad k > 0 \label{eq:arp_acf} \\
P(f) &= \frac{\sigma^{2}}{\left|1 - \sum\limits_{k=1}^p \varphi_{k} e\strut^{-j 2 \pi f k / f_s}\right|^{2}} \label{eq:arp_psd}
\end{align}

AR($p$) models lose the single-parameter interpretation of the process: the timescales. With more than one lag, we must interpret a set of coefficients. Each one measures what its own lag adds beyond the others, so $\varphi_1$ is no longer a correlation and the ACF is no longer a single exponential. However, we will show that AR($p$) model can be factored into a set of processes that correspond to timescales and oscillators.

\subsection{Sources of bias}\label{sec:bias}

Simulations allow the study of timescale estimation bias when analytical solutions do not exist. Timescale estimation is evaluated when the AR(1) assumption fails. Specifically, estimation bias is evaluated under: \textbf{1)} model mismatch, when a process is better described by a higher-order AR($p$) model \cite{broersen_finite_2000}, \textbf{2)} the presence of neural oscillations \cite{berger_uber_1929, buzsaki_neuronal_2004}, \textbf{3)} finite duration or decreasing the number of samples in the signal \cite{broersen_finite_2000, schwarz_estimating_1978, zeraati_flexible_2022}, and \textbf{4)} dynamic, non-stationary timescales associated with behavioral or task states, or neural regions \cite{gao_neuronal_2020, spitmaan_multiple_2020}.

The presence of oscillations is one type of model mismatch. In general, model mismatch is handled by increasing the complexity of the model to AR($p$). Finite bias is a data collection problem that must be addressed at the level of experimental design.

\begin{figure}[htbp]
  \centering
  \includesvg[width=0.95\linewidth]{revision/figures/fig_bias.svg}
  \caption{Timescale estimation bias in the signal (top), autocorrelation (middle) and power spectral density (bottom). The autocorrelations and spectra are theoretical except in the finite-sample column. \textbf{a) Model mismatch.} The constructed AR(7), its exact spectrum and autocorrelation. The AR(1) form does not well describe the ACF or PSD, and higher order AR($p$) model is suggested. \textbf{b) Oscillations.} A damped oscillator in the ACF and a peak in the PSD requires an AR(2) model. \textbf{c) Finite samples.} One second of an AR(1) process. The observed ACF and PSD deviate randomly from their theoretical forms, and the fitted timescale is biased. Longer signals are required to increase confidence in timescale estimates. \textbf{d) Non-stationarity.} A non-constant timescale violates the assumption of stationarity. The AR(1) coefficient increases over the simulation.  No single AR(1) describes all of them, so accurate timescale estimation depends on an appropriate windowing scheme.}
  \label{fig:bias}
\end{figure}

\subsubsection{Model mismatch}

Model mismatch occurs when the AR(1) model is the incorrect model, after other biases have been carefully handled (\figpanel{fig:bias}{a}). This is a generalized extension of oscillatory bias, when there is temporal structure that does not neatly fall into the AR(1) form. In these cases, the AR(1) model is the incorrect model after oscillations have been accounted for, the signal is confirmed to be stationary, and the signal is sufficiently long, yet is not well described by an AR(1) model. When this occurs, the use of higher order AR($p$) offers a solution (eq. \ref{eq:arp}). Increasing the AR order flexibly models PSD and ACF, allowing for forms that deviate from the expected.

The denominator of eq. \ref{eq:arp_psd} is a polynomial: $\varphi_1 (e^{-j 2 \pi f / f_s})^1 + \varphi_2 (e^{-j 2 \pi f / f_s})^2 + \dots + \varphi_p (e^{-j 2 \pi f / f_s})^p$. This allows application of the Weierstrass approximation theorem, suggesting the AR($p$) model can fit any spectral form, i.e. universal approximation, given a sufficiently large order, $p$. The AR($p$) model may overfit other sources of variance, such as finite bias, thus it is desirable to keep $p$ as small as possible. Bayesian information criteria (BIC) may also be used to evaluate a set of models with varying $p$ to determine an optimal, minimized AR order \cite{schwarz_estimating_1978, vidaurre_survey_2013}.

Raising the order has an interpretive cost and it has limited the use of AR($p$) models in this literature. A timescale is a single number, and an AR($p$) fit returns a vector of $p$ coefficients. We show in the Results that an AR($p$) model is interpretable as a set of decaying and damped oscillatory processes, each with a timescale (\figpanel{fig:modes}{}).

\subsubsection{Oscillations}

In the context of neural data, timescales are typically estimated from the autocorrelation of LFPs or binned spike counts \cite{murray_hierarchy_2014}. However both LFPs and single-unit spiking often exhibit rhythmic fluctuations -- oscillations. Oscillatory activity manifests as a cosine in the ACF \cite{brunel_fast_1999} or as a narrow peak in the PSD \cite{he_scale-free_2014, buzsaki_scaling_2013}. An AR(1) model has no term for either, so the timescale it reports absorbs the rhythm and is biased (\figpanel{fig:bias}{b}).

Higher order models capture oscillations. An AR(2) process is a damped oscillator \cite{stralrowski_charts_1970}: a cosine in the ACF under an exponentially decaying envelope, and a peak of finite width in the PSD. Each further pair of coefficients adds another such rhythm, so an AR($p$) model describes the oscillations and the aperiodic decay together, and each rhythm is reported as a frequency and a timescale, the decay of its envelope. Neural rhythms wax and wane, drift in frequency, and often appear as brief bursts \cite{feingold_bursts_2015}, which is what a damped oscillator describes and what the empirical neural recordings show (\figpanel{fig:realdata}{a,c}).

The exception is a perfectly sustained sinusoid, such as line noise or a steady-state response. Its ACF is a cosine that never decays and its PSD is a line at a single frequency, which no AR model produces at any order. Such a rhythm should be excluded rather than modeled. This is possible in the PSD, where it occupies a few bins and can be discarded by robust regression or modeled explicitly as a peak \cite{donoghue_parameterizing_2020}, and not in the ACF, where it perturbs every lag.

\subsubsection{Finite sampling}

Stochastic processes suffer from statistical bias related to finite duration recorded signals (\figpanel{fig:bias}{c}). Decreasing the duration of a time series increases timescale estimation bias, limits the maximum detectable timescale, and decreases confidence. Moving towards unbiased timescale estimation approaches is necessary for studying the windows temporal dynamics of task designs (\figpanel{fig:bias}{d}). Dynamic timescales have been observed in working memory \cite{wasmuht_intrinsic_2018, gao_neuronal_2020}, decision making \cite{ganupuru_flexibility_2019} and auditory processing \cite{chien_constructing_2020} tasks.

In the context of event-related task designs or neuronal spiking activity, bias is often mitigated by averaging the ACFs of trials or binned spike counts \cite{murray_hierarchy_2014}. In these approaches, when there are multiple neurons, channels, or repeat trials, it is recommended to take the median or mean PSD or ACF across recordings, neurons, and/or time periods, prior to model fitting. Averaging across neuronal populations, regions, and time has been shown to reduce estimation bias attributed to short duration recordings \cite{spitzner_mr_2021}. However, averaging limits the ability to examine between-neuron, region, or trial timescale variability. This suggests that across regions or neurons, signals should be stationary and from similar distributions.

\subsubsection{Non-stationarity}

A signal is stationary when its statistics do not change over time, and non-stationary otherwise. In the terms used here, a dynamic signal is one whose AR($p$) parameters, and therefore its timescale, are not constant across the recording. This occurs when a signal spans varying cognitive or neural states. Fit to such a signal, a single model combines distinct processes and reports a timescale that belongs to none of them (\figpanel{fig:bias}{d}). Timescales should therefore be estimated from signals that are homogeneous with respect to anatomical or functional boundaries and to behavioral or task states. This can be done by windowing the recording into periods in which the state is constant, which requires choosing a window size and knowing which windows to combine, or by designing the experiment so that the state is held constant throughout. Windowing carries a cost, since shorter segments increase finite sampling bias and each estimate is less confident, leaving a trade-off between minimizing violations of stationarity and maximizing window size.

\section{Results}\label{sec:results}

\subsection{An AR($p$) model is a set of timescales}\label{sec:modes}

Raising the order removes model mismatch, but an AR($p$) fit returns $p$ coefficients rather than a timescale. The coefficients are not the natural description of the model. Every AR($p$) model is equivalent to a small set of \emph{modes}, and each mode is one way the process rings down after an input: a decay with time constant $\tau$, and possibly a rhythm at frequency $f$ (\figpanel{fig:modes}{a}). The modes are the poles $\lambda_k$ of the model, the roots of its characteristic polynomial (Methods). A real pole is a pure decay and a complex-conjugate pair is a damped oscillation, with
\begin{align}
\tau_k = \frac{-1/f_s}{\log|\lambda_k|}, \qquad f_k = \frac{f_s |\arg \lambda_k|}{2\pi}. \label{eq:mode}
\end{align}
For a real pole, eq. \ref{eq:mode} reduces to eq. \ref{eq:tau}, so the AR(1) model is the special case of a single real pole. A conjugate pair adds a damped cosine to the ACF, $e^{-t/\tau_k}\cos(2\pi f_k t)$, where $f_k$ sets the period and $\tau_k$ sets how fast the ringing dies away; the two are independent. In the ACF the modes add, one decay or damped cosine each (\figpanel{fig:modes}{c}). In the PSD they multiply, one factor each (\figpanel{fig:modes}{b}). Both decompositions are exact (Methods), so an AR($p$) fit can be reported as a short list of $(\tau, f)$ pairs. Not every mode matters equally. Deleting a mode and refitting the overall scale degrades the fitted spectrum by an amount that has a closed form, and we use that degradation as the mode's importance (Methods). The AR(7) of \figpanel{fig:modes}{} has five modes, relaxations at 200, 20 and 2 ms and rhythms at 10 Hz and 40 Hz. Four of them reproduce its spectrum to $R^2 = 0.997$, and the 2 ms relaxation is the one that does not matter.

The modes can be recovered from data, and the coefficients cannot. Across 200 simulations of 60 seconds each, the spectral AR($p$) fit recovered every mode of the AR(7), with median errors in $\tau$ between 1.9\% and 8.0\% and both rhythm frequencies to within 0.6\%. The first coefficient of the same fits was $5.92 \pm 0.01$. It is not a correlation and has no timescale. An AR(1) fit to the same signals ran to the 20 s bound of its search in every simulation, which is the fit reporting that no single decay describes the spectrum.

\begin{figure}[htbp]
  \centering
  \includesvg[width=\linewidth]{revision/figures/fig_modes.svg}
  \caption{An AR($p$) model read as modes. The constructed AR(7) has five modes: relaxations at 200, 20 and 2 ms, and rhythms at 10 Hz (150 ms envelope) and 40 Hz (60 ms envelope). A real pole costs one order and a conjugate pair costs two, giving $3 + 4 = 7$. \textbf{a) Modes.} The response to a single input is the sum of the modes, ordered slowest to fastest and drawn at unit amplitude. Blue is a relaxation mode, red an oscillatory mode, black the summed response. \textbf{b) Spectrum.} The modes multiply, one factor each; the grey curve is their product, the spectrum of the process. \textbf{c) Autocorrelation.} The modes add, as exponential decays for real poles and damped cosines for conjugate pairs; the grey curve is their sum. The weights are set by the poles: the 200 ms relaxation enters with weight 0.90, the 10 Hz rhythm 0.13, the 20 ms relaxation $-0.03$ and the 40 Hz rhythm $2 \times 10^{-5}$, so the 40 Hz mode, prominent on the logarithmic axis of (b), is invisible in (c).}
  \label{fig:modes}
\end{figure}

\subsection{Recovering timescales in simulation}\label{sec:sim}

We tested the six estimators, three AR(1) and three AR($p$), on processes whose timescale is known: an AR(1) process, an AR(2) oscillator, and the AR(7) of \figpanel{fig:modes}{}, each built with a 200 ms timescale (\figpanel{fig:benchmark}{}; Methods). When the process is AR(1), all six estimators recover the timescale to within 5\%, and BIC selects $p = 1$. When it is not, the AR(1) estimators disagree by up to a factor of 112, and every one of the disagreeing numbers is reproducible across simulations. The order-selected AR($p$) fits recover the oscillator's 200 ms envelope decay to within 1\% and its 10 Hz rhythm to within 0.01 Hz, and they recover the AR(7)'s dominant relaxation to within 9\%, at 183 to 195 ms. The oscillator result holds across a decade of envelope decays and three octaves of frequency (Supplementary Results). A rhythm that is part of the dynamics should therefore be modeled, and modeling it returns its frequency and its envelope decay. A sinusoid added to the signal is different. It is not autoregressive at any order, so no order repairs the ACF fit, while robust regression in the spectrum removes it (Supplementary \figpanel{fig:s_grid}{}).

A short recording biases a timescale and makes it imprecise, and both effects depend on the recording only through the ratio $T/\tau$ of its duration to the timescale (\figpanel{fig:finite}{}). The relative bias is about $-2\tau/T$ and the relative standard error about $\sqrt{2\tau/T}$ (Methods), so a bias under 10\% needs about $20\tau$ of data and a standard error under 10\% needs about $200\tau$, at any sampling rate. For a 100 ms cortical timescale, that is 2 and 20 seconds of stationary data per estimate. The sampling distribution is skewed to the right (\figpanel{fig:finite}{c}), so a short recording more often underestimates the timescale but errs further when it overestimates. On four seconds of the 200 ms AR(1) process, twenty timescales, the central 80\% of AR($p$) estimates spans a factor of five (\figpanel{fig:benchmark}{d}, last column). This is what makes non-stationarity the hardest bias to handle. A window short enough to be stationary is exactly the one that biases the estimate, and the better response is to design experiments whose states last tens to hundreds of timescales.

\begin{figure}[htbp]
  \centering
  \includesvg[width=\linewidth]{revision/figures/fig_benchmark.svg}
  \caption{The AR($p$) estimators on four simulated cases. Columns: an AR(1) process, the AR(2) oscillator of \figpanel{fig:bias}{b}, the AR(7) of \figpanel{fig:modes}{}, and four seconds of the same AR(1), a finite sample of twenty timescales. Every column has the same constructed timescale, 200 ms: the AR(1)'s $\tau$, the envelope decay of the oscillator, and the dominant relaxation of the AR(7). The first three are recorded for 100 s at 500 Hz, 200 repetitions. \textbf{a) Signal.} 1.5 s of each. \textbf{b) Spectrum.} One simulation (black) with the order-selected AR($p$) spectral fit (blue, dashed). \textbf{c) Autocorrelation.} Over the analysis band (black) with the order-selected AR($p$) lag-domain fit (blue, dashed). \textbf{d) Estimates.} The timescale each AR($p$) fit reports, for the fit to the ACF and to the PSD: box is the interquartile range over repetitions, line the median, whiskers the 10th to 90th percentile, on one axis for the row. The dashed line is the constructed timescale, red for the oscillator, whose timescale is the envelope decay of its rhythm.}
  \label{fig:benchmark}
\end{figure}

\begin{figure}[htbp]
  \centering
  \includesvg[width=\linewidth]{revision/figures/fig_finite.svg}
  \caption{Finite samples. \textbf{a, b) Confidence intervals.} The intervals around the estimated AR(1) coefficient narrow as the number of samples increases from 100 (a) to 1000 (b), and widen as $\varphi$ approaches one, which is where neural timescales lie. Estimates and intervals (gray) were averaged over 1000 simulations. \textbf{c) Sampling distribution.} The density of $\hat\tau/\tau$ for the time-domain AR(1) fit at recording lengths of $10\tau$, $100\tau$ and $1000\tau$ (150 simulations each). The error is shown as a ratio because bias and spread both scale with $\tau$ (eq. \ref{eq:finite}), so one panel covers every timescale. The short recordings are centred below one, by $-2\tau/T$: $-20\%$ at $T = 10\tau$, $-2\%$ at $100\tau$ and $-0.2\%$ at $1000\tau$. The sample size that matters is $T/\tau$, the number of independent stretches, not the number of samples.}
  \label{fig:finite}
\end{figure}

\subsection{Timescales in neural recordings}\label{sec:realdata}

We applied one pipeline to eight recordings spanning fMRI, MEG, EEG, ECoG, LFP and spiking activity in humans, macaque, mouse and rat (\figpanel{fig:realdata}{}; Methods). The estimates reproduce the separation of scales that motivates the hierarchy literature. Field-potential timescales were tens of milliseconds, 10 to 48 ms by the AR(1) spectral fit; spike counts gave 198 ms; and precuneus BOLD gave 2.9 s, between 1.0 and 9.3 s across participants.

No recording was AR(1). Every field potential contained a narrowband peak, in six of the eight the ACF oscillates rather than decaying, and only the AR($p$) model follows it (\figpanel{fig:realdata}{a,c}). An AR(1) model left a median of 13\% of the log spectrum unexplained, between 5\% and 52\% across recordings, and an order-selected AR($p$) model left 1.5\%, between 0.2\% and 38\%. Because no recording satisfied the assumption, the estimators disagreed. The ratio of the largest to the smallest reported timescale ranged from 1.3 to 11 across the ACF and PSD fits, with a median of 5.4, and from 1.6 to 115 with the time-domain fits included, with a median of 11. The largest disagreements were in the recordings with the strongest rhythms. There the AR(1) fits report a fast decay, and the AR($p$) fits report a slower relaxation once the rhythm is described rather than absorbed.

The modes say what each recording contains (\figpanel{fig:realdata}{b,d,e}). Keeping the modes whose timescale the analysis band can resolve (Methods) leaves between one and eleven per trace, and the figure shows the three most important. In resting-state fMRI one mode reproduces the fitted spectrum, a relaxation of 3.4 s, and in the spike counts two, a 420 ms relaxation and a slower rhythm near 1.5 Hz in half the units. Every field-potential recording needs seven or more, and the second is not a detail. In posterior alpha, the components are a 61 ms rhythm at 10 Hz, a 33 ms rhythm at 22 Hz and a 70 ms relaxation; in rat hippocampus, an 82 ms rhythm at 9 Hz, a 46 ms rhythm at 21 Hz and a 71 ms relaxation; in sensorimotor MEG, a 26 ms rhythm at 15 Hz, a 39 ms relaxation and a 19 ms rhythm at 27 Hz; in macaque ECoG, a 34 ms rhythm at 15 Hz beside a 126 ms relaxation. These components are consistent across the traces of a recording. The dominant component's timescale varied between the 10th and 90th percentiles of traces by a factor of 1.8 to 8 across channels and participants, with a median of 3.3 (\figpanel{fig:realdata}{e}), against disagreements of up to 115-fold between estimators on the same traces. A single reported $\tau$ is one of these components, and which one depends on the estimator. A higher-order model identifies the dominant relaxation mode and reports the rhythms beside it. Cutting the recordings into windows shows the remaining bias. The spread of estimates across windows exceeded the AR(1) sampling bound in every recording, by a median factor of about 13 (Supplementary \figpanel{fig:s_windows}{}). Part of that spread is variation between channels and participants and part is change over time, and sampling explains none of it, so the timescale of a long recording is an average over a quantity that varied within it.

\begin{figure}[htbp]
  \centering
  \includesvg[width=\linewidth]{revision/figures/fig_realdata.svg}
  \caption{Timescales in eight neural recordings, six modalities and four species, each analysed with the identical pipeline; every panel pools all traces of a recording, and every panel describes one model, the order-selected AR($p$) fit to the spectrum. \textbf{a) PSD fit.} Median spectrum across traces (black) with the median AR($p$) fit (blue, dashed). Rows are labelled by species and modality. \textbf{b) PSD components.} The factors the fitted spectrum is a product of, for the three most important modes the band can resolve, each drawn relative to its own geometric mean and coloured by importance rank; a relaxation is a shoulder and a rhythm a peak. \textbf{c) ACF.} Autocorrelation over the analysis band (black) with the autocorrelation implied by the fit in (a) (blue, dashed), both derived from the band-limited spectrum in the same way (Methods). The lag axis is cropped to the decay. \textbf{d) ACF components.} The same three modes as additive components of the autocorrelation, which add exactly to the dashed curve in (c); a relaxation is an exponential decay and a rhythm a damped cosine. \textbf{e) Timescales.} The timescale of each of those components over the traces that resolve them, in the same colours: box is the interquartile range, line the median, whiskers the 10th to 90th percentile. A row with one box is a recording that one timescale describes, and a row with three is one that needs three.}
  \label{fig:realdata}
\end{figure}

\section{Discussion}\label{sec:discussion}

A neural timescale is a parameter of a model. When the model is the AR(1) process, as it is for an exponential fit to the ACF, a Lorentzian fit to the PSD, or a lagged regression on the signal, $\tau$ is the relaxation time of that process, and reporting it as a property of a region, a state or a subject is reasonable. When the model is wrong, $\tau$ is still precise, still reproducible and still comparable across conditions, but it is no longer the quantity the interpretation is about. Our simulations show that the two situations look identical from the outside, and our recordings show that neural data are usually in the second one. Three checks tell them apart. The first is agreement between independent estimators, which is within a few percent when the data are AR(1) and up to 112-fold in simulation and 115-fold in the recordings when they are not. The second is the fraction of the spectrum or ACF that the fitted model leaves unexplained. The third is the order BIC selects. None of the eight recordings passed all three.

The response to model mismatch is a higher order, and the result is readable. An AR($p$) model is a list of decay times and rhythms (\figpanel{fig:modes}{}), the dominant relaxation among them is the timescale an AR(1) fit was trying to report, and the rhythms are reported beside it with their own decay. In every field-potential recording, that list has more than one entry that matters (\figpanel{fig:realdata}{}), and the entries are consistent across channels and participants where the estimators are not. An oscillation that is part of the dynamics belongs on that list; an oscillation added to the signal, such as line noise or a steady-state response, is not autoregressive at any order and must be removed instead, which the spectrum allows and the ACF does not. Finite samples set a floor of about $200\tau$ per estimate, and non-stationarity pushes against it: a window short enough to be stationary is often too short to be precise. The spread across windows of the recordings was about an order of magnitude larger than finite sampling predicts, so a confidence interval from the sampling formula understates the uncertainty of a windowed estimate, and one from the spread across windows does not. These biases rarely occur alone. Every recording here had a rhythm, a finite duration and a timescale that drifted, and they have to be handled together.

\subsection{Recommendations}

We suggest the following when reporting a neural timescale. \textit{First}, state the model rather than only the estimator, for example ``the relaxation time of an AR(1) fit to the 1-100 Hz spectrum'' rather than ``the neural timescale''. \textit{Second}, report the ratio of recording length to estimated timescale, and treat estimates below roughly $200\tau$ as precision-limited. \textit{Third}, report the fraction of the spectrum or autocorrelation left unexplained by the fitted model. \textit{Fourth}, report at least one independent estimator; for recordings longer than $200\tau$, disagreement beyond about 1.5-fold is evidence of model mismatch rather than of sampling (Methods). \textit{Fifth}, if the data contain a strong narrowband rhythm, prefer spectral-domain estimation and state whether the rhythm is being modeled or excluded. \textit{Sixth}, if an AR($p$) model is used, report the modes as $(\tau, f)$ pairs with their importances and the number needed to reproduce the fit, not the coefficient vector. \textit{Finally}, when windowing, take the confidence interval from the spread across windows rather than from the sampling formula.

\subsection{Limitations}

Our simulations treat stationary processes, and non-stationarity is addressed only through windowing. It remains the hardest of the four biases, because shortening windows to reduce it increases finite sampling bias. Seven of the eight recordings are field potentials or a hemodynamic signal, and only one is a spike train. Binning and Poisson variability add a flat floor to a spike train that the AR spectral model cannot represent \cite{zeraati_flexible_2022, meyer_inferring_2019}; we handle this only by restricting the fit band, and a model with an explicit observation-noise term would treat it properly. The real-data analysis applies one pipeline to every recording, so the preprocessing a specialist would apply, artifact rejection, burst detection and oscillation removal, is absent by design, and the mismatch we report is an upper bound on what careful per-dataset treatment would leave. The band-limited autocorrelation we fit (Methods) is not the raw autocorrelation of the published literature, and a timescale from one is not directly comparable to a timescale from the other; on one recording the two differ by a factor of sixty. Finally, the mode decomposition is exact for the fitted model but inherits its assumptions. It is a parsimonious description of a spectrum, and whether a given mode corresponds to a separable biological process is a question the model cannot answer. For example, the sum of two independent AR(1) processes is not itself autoregressive, and its fitted poles are pulled away from both original timescales.

\section{Methods}\label{sec:methods}

\subsection{Estimators}

Six estimators were compared. They differ in the domain they are fit in and in the model order, and they share everything else. Each uses the same mean-centered signal and the same Welch power spectrum \cite{welch_use_1967}, so differences between them cannot come from preprocessing. In the time domain, the AR(1) estimator is a least-squares regression of $x_t$ on $x_{t-1}$, and the AR($p$) estimator uses Burg's method \cite{burg_maximum_1975}. In the autocorrelation domain, the model is $A\,\rho_p(k) + B$, where $\rho_p$ is the exact AR($p$) autocorrelation, $A$ absorbs a white-noise floor at lag zero, and $B$ is a constant offset. For $p=1$ this is the exponential-plus-offset model standard in this literature \cite{murray_hierarchy_2014}. In the spectral domain, the model is eq. \ref{eq:arp_psd} with a free scale, fit to $\log_{10}$ power, as in spectral parameterization methods \cite{donoghue_parameterizing_2020}.

A spectral fit is restricted to an analysis band, 1-100 Hz for most of the electrophysiological recordings here, whereas the sample autocorrelation of the raw signal contains everything in the recording, including drift below that band. Fitting one estimator to the band and the other to the whole signal would compare two answers to two different questions. We therefore compute the autocorrelation from the same band-limited spectrum the spectral estimators see, by inverse Fourier transform. Bins excluded as line noise are filled by log-linear interpolation from their neighbours rather than zeroed, so that notching does not ring in the lag domain. The difference is large: on rat hippocampal CA1/CA3, the raw autocorrelation gives an AR(1) timescale of 938 ms, set by slow drift the spectral fit never saw, against 14 ms once both estimators describe the same band. Welch's estimate is the average periodogram of tapered segments, so the autocovariance it implies is the true autocovariance multiplied by the taper's own autocorrelation, which for a Hann window has fallen to 0.66 at a lag one quarter of the segment length. We divide this taper out and discard the lags at which it retains less than a fifth of its energy. A model fit to the spectrum is carried into the lag domain by the same route, its spectrum evaluated on the same grid, restricted to the same band and inverse transformed (\figpanel{fig:realdata}{c}). The components of \figpanel{fig:realdata}{d} are the partial fractions of the fitted spectrum, one term per pole, transformed the same way, and they add exactly to the model's band-limited autocorrelation.

Robust (soft-L1) regression is used in the spectral domain only, and only when the signal contains a rhythm: in simulation, for the AR(2) oscillator and the added sinusoid; in real data, for every recording. A rhythm is an outlier in the spectrum, occupying a handful of bins, and can be down-weighted there. In the lag domain it perturbs every lag, so down-weighting large residuals discards the fit rather than the rhythm. Robustness trades one thing a fit can report for another. On an AR(1) background at $\tau = 50$ ms plus a broadband AR(2) rhythm at 10 Hz, the robust spectral fit recovered the aperiodic timescale at 53 ms over 20 simulations while discounting the rhythm; without robustness the fit found the rhythm at a median of 9.4 Hz and returned the aperiodic timescale at 68 ms. A study that wants both should report them from two fits.

Both the autocorrelation and spectral fits are optimized over reflection coefficients rather than over the AR coefficients directly. Every point in the box $|k_i| < 1$ is a stationary model, the parameters are all of order one, and the fit cannot wander into a region where mode timescales are undefined. Without this reparameterization, 48\% of the spectral fits to real data returned non-stationary parameters. Multiple starting points are used because the objective is non-convex when a rhythm is present. Two further requirements are placed on the fitted model, because it is being used to measure a timescale. First, the model must contain a relaxation mode. A converged, well-fitting AR($p$) can place every pole in a conjugate pair, in which case it has no decay time to report, and before the constraint a third of the fits to our recordings did exactly that. We therefore factor one real pole out of the characteristic polynomial, $A(z^{-1}) = (1 - a z^{-1})A_{\mathrm{rest}}(z^{-1})$ with $0 < a < 1$, and parameterize the remainder by reflection coefficients. The requirement costs one order and is applied at every order, so on an AR(2) oscillator BIC selects order three; it changes each estimator's own objective by a median of $+0.003$ (lag domain) and $-0.000$ (spectrum) across our recordings. Second, the reported aperiodic timescale is the most important non-rhythmic mode, where non-rhythmic means $f\tau < 1$, a mode that rings down before completing a cycle; a real pole qualifies trivially. The guaranteed real pole is guaranteed to exist, not to matter: on mouse LFP the constrained spectral fit parks its required pole at 0.6 ms with negligible importance, while the aperiodic structure sits in an 88 ms pair ringing 0.57 of a cycle. Importance, not pole type, decides.

Spectra are estimated by Welch's method with 4 s windows and 50\% overlap; for fMRI, sampled every 2 s, the window is 128 s, and for the spike counts, binned at 50 ms, it is 60 s. Models fit in either domain are scored in both: spectral fit is the coefficient of determination of $\log_{10}$ power over the analysis band, with the overall scale refit, and autocorrelation fit is the coefficient of determination over the fitted lags, with $A$ and $B$ refit.

\subsection{Model order}

Order is selected by BIC \cite{schwarz_estimating_1978} in the domain the model is fit in: from the conditional Gaussian likelihood in the time domain, and from the residual sum of squares in the autocorrelation and spectral domains. The search runs over a log-spaced grid of orders up to 20. In the lag domain, neighbouring lags are strongly dependent, so the criterion uses an effective sample size, the number of lags divided by the lag at which the autocorrelation first falls below $1/e$. When the selected order reaches the top of the range, this is recorded and reported, because it indicates a process the AR class cannot represent at any practical order.

\subsection{Modes and importance}

The AR($p$) recursion in eq. \ref{eq:arp} has a characteristic polynomial $z^p - \varphi_1 z^{p-1} - \cdots - \varphi_p$, whose roots $\lambda_k$ are the poles. Stationarity requires $|\lambda_k| < 1$. Dividing by $z^p$ gives
\begin{align}
1 - \sum_{k=1}^{p} \varphi_k z^{-k} = \prod_{k=1}^{p} (1 - \lambda_k z^{-1}), \label{eq:factor}
\end{align}
so the spectral denominator in eq. \ref{eq:arp_psd} factorizes into first-order terms, and each pole contributes a decay time and, when complex, a frequency (eq. \ref{eq:mode}). Because the coefficients are real, complex poles come in conjugate pairs, and we combine each pair into a single damped oscillatory mode, equivalent to a real AR(2) factor. A real pole is a pure relaxation mode. The term comes from linear systems theory, where the modes of a linear system are the eigenvalues that govern how it relaxes; the decomposition is a general result for linear time series \cite{west_time_1997} that has been used to track the components of a fitted autoregression through EEG seizures \cite{west_evaluation_1999}.

The factorization is of the transfer function, not of the signal into independent AR(1) processes: the factors multiply in the spectrum, and an AR($p$) process is a cascade of first-order filters driven by one input. The same poles appear additively in the time domain. The impulse response is $\psi_n = \sum_k c_k \lambda_k^n$, and the autocovariance decomposes the same way,
\begin{align}
\gamma(k) = \sum_j A_j \lambda_j^{\,k}, \label{eq:acfmodes}
\end{align}
so, grouping conjugate terms, the autocorrelation of any AR($p$) process is a sum of exponential decays and damped cosines, one per mode. The product form and the sum form are the factored and partial-fraction representations of the same rational function. Each residue $A_j$ depends on every pole, so a mode's weight in the autocorrelation is not a free parameter, which is why a lag-domain AR($p$) fit to an added sinusoid is biased (Supplementary \figpanel{fig:s_grid}{}), and why the process variance does not divide into per-mode shares. In our implementation both decompositions reproduce the model to better than one part in $10^7$.

Writing $g_k(f)$ for the log of the $k$th pole factor, the model log spectrum is $\log \sigma^2 - \sum_k g_k(f)$, so modes are additive in log power. Deleting mode $k$ and refitting the overall scale by least squares leaves a residual whose root-mean-square in $\log_{10}$ power is exactly $\mathrm{sd}_f\{g_k\}/\ln 10$, and we use this as the mode's importance. An Itakura-Saito variant, equal to the asymptotic Whittle log-likelihood degradation \cite{whittle_estimation_1953, itakura_analysis_1968}, gives the same ranking in every case we examined. Importance is measured on the log spectrum and is not a share of the variance: the 40 Hz mode of \figpanel{fig:modes}{} ranks second in importance and carries $2 \times 10^{-5}$ of the variance. The error of a model that keeps a subset of modes is a quadratic form in the kept set, so the best subset of each size can be evaluated in closed form, and the dominant modes are the smallest subset that reproduces the fitted log spectrum with $R^2 \ge 0.99$. The count depends on this cutoff: at $R^2 \ge 0.95$ the field potentials need 6 to 11 modes and at $0.99$ they need 7 to 11, while at $0.90$ they need 1 to 6. We report the strict cutoff, and the three modes shown in \figpanel{fig:realdata}{} are the most important under any of them.

For the real recordings, a mode is reported only if the analysis band can resolve its timescale. A mode slower than $1/(2\pi f_{\mathrm{lo}})$, where $f_{\mathrm{lo}}$ is the lower edge of the fitted band, has its knee below the band, so its timescale is an extrapolation of the model past the data. Left in, these near-unit-root poles report timescales longer than the recording and carry high importance precisely because deleting them changes the lowest frequencies. The criterion discards 7\% of the fitted spectral modes. Likewise a rhythmic mode counts as detected only if its frequency lies inside the fitted range.

\subsection{Simulated processes}

Four processes were used in the benchmark, each with a closed-form spectrum and autocorrelation, simulated at 500 Hz for 100 s with 200 repetitions. Each isolates one way the AR(1) assumption can fail, and every process with a timescale has the same one, 200 ms, so that the estimates can be read on one axis. \textit{AR(1)} with $\tau = 200$ ms is the positive control. \textit{AR(2) oscillator} has poles at $|\lambda| e^{\pm i\theta}$ with $|\lambda| = e^{-1/(\tau f_s)}$ and $\theta = 2\pi f_0/f_s$, here $f_0 = 10$ Hz and $\tau = 200$ ms, two cycles. \textit{AR($p$)} is the AR(7) of \figpanel{fig:modes}{}: relaxations of 200, 20 and 2 ms, and rhythms at 10 Hz ringing for 150 ms and at 40 Hz for 60 ms. It is the only process inside the model family but above first order, so it isolates order mismatch from model mismatch; its two slow relaxations have comparable importance an order of magnitude apart, and its rhythms ring for more than a cycle, so the most important non-rhythmic mode is unambiguously the 200 ms relaxation. \textit{Power law} is a long-memory ARFIMA process \cite{granger_introduction_1980, hosking_fractional_1981}, simulated exactly by circulant embedding \cite{davies_tests_1987}; it has no characteristic timescale, so there is nothing to recover (Supplementary Results). The finite-sample column of \figpanel{fig:benchmark}{} is the same AR(1) recorded for four seconds, with the Welch window and the autocorrelation a quarter of the recording, as in the windowing analysis. An alternative higher-order process, an AR(5) with a 250 ms relaxation, a 5 Hz mode ringing half a cycle and a 50 Hz rhythm, is simulated alongside and can be drawn in place of the AR(7) in every figure.

Two sweeps extend the benchmark. The damped-oscillator sweep takes the AR(2) oscillator over envelope decays of 20, 50 and 150 ms and frequencies of 5, 10, 20 and 40 Hz. The added-sinusoid sweep adds $A\cos(2\pi f_0 t + \phi)$ with random phase to an AR(1) with $\tau = 50$ ms, over six frequencies from 4 to 40 Hz and six power ratios from one eighth to four (Supplementary \figpanel{fig:s_grid}{}). Both use 20 repetitions per cell. Outside AR(1), the target has to be stated, and we take it to be whatever the constructed process gives under the same rule the estimator applies, computed from the true poles: the relaxation time for AR(1), the envelope decay for the oscillator, and the most important non-rhythmic mode for the AR(7). Every theoretical spectrum and autocorrelation was verified against long simulations.

\subsection{Finite-sample theory}

For a mean-corrected least-squares AR(1) fit on $n$ samples, the classical results are $\mathbb{E}[\hat\varphi] - \varphi \approx -(1+3\varphi)/n$ \cite{marriott_bias_1954, kendall_note_1954} and
\begin{align}
\mathbb{E}[ (\hat{\varphi} - \varphi)^2 ] &= \frac{1 - \varphi^2}{n-1}. \label{eq:var}
\end{align}
Propagating both through eq. \ref{eq:tau} to second order, and writing $T = n/f_s$ for the recording duration, gives, in the regime $\tau f_s \gg 1$ that all neural recordings occupy,
\begin{align}
\frac{\mathbb{E}[\hat\tau] - \tau}{\tau} \approx -\frac{2\tau}{T}, \qquad
\frac{\mathrm{sd}[\hat\tau]}{\tau} \approx \sqrt{\frac{2\tau}{T}}. \label{eq:finite}
\end{align}
Both depend on the recording only through $T/\tau$, and both were verified against Monte Carlo. This is why estimation error is reported throughout as the ratio $\hat\tau/\tau$: it has the same distribution for a 10 ms and a 1 s timescale at the same $T/\tau$, whereas absolute and squared errors scale with $\tau$ and $\tau^2$. Confidence intervals are obtained by mapping the interval for $\hat\varphi$ through eq. \ref{eq:tau}, which preserves coverage because the map is monotone. The intervals of \figpanel{fig:finite}{a,b} are the least-squares intervals of eq. \ref{eq:var}, evaluated at each estimate and averaged over simulations; the weighted least-squares/Cauchy construction \cite{gallagher_small_2008} refines them for $\varphi$ very close to one.

\subsection{Recordings and processing}

Eight publicly available recordings were used, selected to span species, recording modality and oscillatory regime. \textit{fMRI}: ABIDE resting-state BOLD \cite{di_martino_autism_2014}, twenty participants from the UM site, preprocessed by C-PAC without temporal filtering, TR 2 s, precuneus left and right. The unfiltered pipeline is used deliberately, since the usual 0.01-0.1 Hz band-pass removes most of the spectrum a timescale is estimated from, and one named region is used rather than an average over regions, because regions differ in exactly the property being measured. \textit{MEG}: the FieldTrip cortico-muscular coherence recording \cite{schoffelen_neuronal_2005, oostenveld_fieldtrip_2011}, the 24 sensorimotor magnetometers with the strongest beta rhythm, during isometric contraction. \textit{EEG}: MPI-LEMON eyes-closed rest \cite{babayan_mind_2019}, four posterior channels in each of five participants. \textit{ECoG}: primary motor cortex in 23 people with Parkinson disease before DBS \cite{cole_nonsinusoidal_2017}, and 24 channels of macaque visual cortex during a grating task \cite{nagasaka_multidimensional_2011}. \textit{LFP}: 24 channels of mouse Neuropixels in VISam during natural movie viewing \cite{siegle_survey_2021}, and 24 channels of rat hippocampal CA1/CA3 \cite{mizuseki_theta_2009, mizuseki_hc3_2013}. \textit{Spiking}: spike counts of 24 VISp units from the same Neuropixels session \cite{siegle_survey_2021}, binned at 50 ms over the same 300 s of natural movie.

Electrophysiological recordings were resampled to 250 Hz and fit over 1-100 Hz, with line-noise fundamentals and harmonics excluded from the spectral fits; the MEG is fit over 1-95 Hz because its spectrum carries a sharp artefactual peak at 100 Hz. fMRI is fit over 0.01-0.24 Hz. The spike counts are fit over 0.05-2 Hz rather than the 0.05-8 Hz their sampling rate allows, because a spike train contributes a flat Poisson floor that the AR spectral model has no term for, and above about 1.5 Hz the binned-count spectrum is within a factor of 1.5 of that floor. Fitted over the full band, the AR(1) model was fitting mostly shot noise and left 95\% of the log spectrum unexplained; the restricted band improves this tenfold. No other filtering, artifact rejection or oscillation removal was applied. Halving or doubling each recording's spectral window changed the median spectral AR($p$) timescale by less than 5\% in motor ECoG, mouse LFP and rat LFP, by 6-18\% in MEG, macaque ECoG, fMRI and the spike counts, and by about 50\% in the EEG, whose traces are the shortest at 60 s.

\subsection{Windowing}

Each recording was cut into non-overlapping windows of increasing length, from 0.5 s upward, and every window was analysed with the identical pipeline, with the Welch window and the autocorrelation a quarter of the window (Supplementary \figpanel{fig:s_windows}{}). Windows are pooled across the traces of a recording, so the spread across windows includes variation between channels and participants as well as over time. Window length is measured in units of the timescale being estimated rather than in seconds, since a 1 s window is generous for a 4 ms timescale and far too short for a 1 s one. The spread of estimates across windows of a given length is compared with the AR(1) sampling bound of eq. \ref{eq:finite} at that length.

\subsection{Thresholds}

Three thresholds are used in the text. The $200\tau$ floor is eq. \ref{eq:finite}: at that length the relative standard error of a single estimate is 10\%. The 1.5-fold rule for disagreement between estimators is set by the same equation: two independent estimates at $T = 200\tau$ each have a 10\% standard error, so their ratio exceeds 1.5 in well under 1\% of cases from sampling alone, and in our AR(1) simulations at $T = 500\tau$ the order-selected ACF and PSD fits disagreed by more than 1.5-fold in 4.5\% of repetitions and the AR(1) fits never did. Below $200\tau$ the rule does not apply: at $T = 20\tau$, sampling alone made the two AR($p$) fits disagree by more than 1.5-fold in half of the repetitions. The $f\tau < 1$ rule for a non-rhythmic mode is definitional: a mode that does not complete one cycle before its envelope has decayed is a damped relaxation, not an oscillation.

\subsection{Code and data availability}

All analyses use an open-source Python module with 200 unit tests, covering exact reconstruction of the AR($p$) spectrum and autocorrelation from the mode decomposition, agreement between the closed-form mode importance and brute-force deletion, agreement of every theoretical spectrum and autocorrelation with long simulations, the guaranteed relaxation mode, and agreement of eq. \ref{eq:finite} with Monte Carlo. Every figure is drawn by its own script from cached result tables, and the simulation stages, tables and figures are tracked as one reproducible pipeline. Code, figure scripts, result tables and the extracted traces are available at \href{https://github.com/voytekresearch/timescale-methods}{https://github.com/voytekresearch/timescale-methods}. All recordings are public: ABIDE, MPI-LEMON, the MNE/FieldTrip CMC dataset, Cole et al. 2017, NeuroTycho, the Allen Brain Observatory and CRCNS hc-3.

\backmatter

\clearpage
\section*{Supplementary Information}\label{sec:supp}

\setcounter{figure}{0}
\renewcommand{\thefigure}{S\arabic{figure}}
\setcounter{table}{0}
\renewcommand{\thetable}{S\arabic{table}}
\setcounter{equation}{0}
\renewcommand{\theequation}{S\arabic{equation}}

\subsection*{Supplementary Results}\label{sec:sresults}

\subsubsection*{Benchmark in detail}

For the AR(2) oscillator, the three AR(1) estimators return 10, 139 and 251 ms for a 200 ms envelope decay, while all three AR($p$) fits recover it to within 1.1\% and the rhythm to within 0.01 Hz. Across the damped-oscillator sweep, a decade of envelope decays and three octaves of frequency, the AR($p$) spectral fit stays within 3\% of the envelope decay and the AR($p$) ACF fit within 9\%, while the AR(1) fits are wrong by factors between 0.016 and 133 that depend on both the decay and the frequency. For the AR(7), the three AR($p$) estimators return 183 to 195 ms against a target of 200 ms, while the three AR(1) estimators spread over a factor of 112: 179 ms for the autocorrelation fit, which the 200 ms relaxation happens to dominate, 1.57 s for the lagged regression, and the 20 s bound of its search for the spectral fit. The long-memory process produces estimates from 4 to 86 ms, a twenty-fold range, every one of them stable across repetitions, and the lag-domain fit returns an estimate proportional to the recording length, $\hat\tau \propto T$ over three decades, where the same estimator on an AR(1) process is flat.

Model mismatch tracks these failures. An AR(1) spectral fit leaves 0.7\% of the log spectrum unexplained for the AR(1) process, but 27\% for the AR(2) oscillator, 67\% for the AR(7) and 26\% for the power law; AR($p$) reduces these to between 0.03\% and 4.0\%. BIC returns $p=1$ for AR(1), $p=3$ for the oscillator (its order plus the required real pole), $p=8$ for the AR(7), and runs to the top of the search range for the power law. Fitting each process's exact autocorrelation, which is what any consistent estimator converges to with infinite data, the mismatch collapses to zero at the correct order for the three autoregressive processes and only shrinks for the power law, which no finite order removes.

The domain a model is fit in changes the model. Across the three processes that are not AR(1), fits to the spectrum describe the autocorrelation nearly as well as fits to the autocorrelation do, leaving between 0.2\% and 7\% of it unexplained. Fits to the autocorrelation describe the spectrum badly, leaving 15\% unexplained for the oscillator, 11\% for the AR(7), and more than the total spectral variance for the long-memory process. Both domains use the same model, reparameterization, order selection and band, so this is not an optimization artifact. A model estimated in the frequency domain describes the signal in either domain; a model estimated in the lag domain does not.

\subsubsection*{Added sinusoids}

A sinusoid added to an AR(1) process is not autoregressive at any order: its spectrum contains a line that no product of poles reproduces, and its ACF contains a cosine that never decays (Supplementary \figpanel{fig:s_grid}{}). As the rhythm grows from one eighth to four times the aperiodic power, the AR(1) ACF fit falls from the true 50 ms to as little as 3 ms, and raising the order changes the failure without removing it: over the same grid, the AR($p$) ACF fit ranges from half to twelve times the truth. The weight of each term in an AR($p$) autocorrelation is set by the poles, and a mixture of an exponential and an undamped cosine has a mixing weight that no set of poles can reproduce. In the spectrum the rhythm occupies a few bins, robust regression discards them, and both spectral fits stay within 3\% of the truth at every frequency and power tested.

\subsubsection*{Windowing in detail}

Measured in units of the timescale, the median spectral AR($p$) estimate settled to within 20\% of its longest-window value once the window exceeded between 40 and 2300 timescales, depending on the recording, with a median of about 550; at the shortest windows it sat between 13\% and 82\% of that value. Taking the autocorrelation from the raw signal instead of from the analysis band makes this far worse: estimates then collapse by up to two orders of magnitude at the shortest windows and climb steadily with window length, which reads as a window-length effect on the brain's timescales but is drift below the fit band entering the lag domain. At the longest windows available, every estimator exceeded the sampling bound of eq. \ref{eq:finite} in all eight recordings, by a median factor of between 11 and 17 and by up to 44, and this does not go away when the autocorrelation is corrected.

\subsubsection*{The integrated timescale}

As a model-free reference we use $\tau_{\mathrm{int}} = \int_0^\infty \rho(t)\,dt$ \cite{sokal_monte_1997, priestley_spectral_1981}, evaluated with the trapezoid convention so that a sampled exponential returns its continuous-time $\tau$ rather than $\tau + 1/(2f_s)$. Equivalently, $\tau_{\mathrm{int}} = S(0)/(4\gamma(0))$. It equals $\tau$ for an AR(1) process, is the variance-weighted mean of component timescales for a sum of independent AR(1) processes, and diverges for a long-memory process. It is not the envelope decay time of an oscillator: the anticorrelated troughs of a damped oscillator's autocorrelation cancel most of the integral, which is why the AR(2) oscillator of the benchmark has $\tau_{\mathrm{int}} = 2.5$ ms and an envelope decay of 200 ms. How long activity stays correlated and how fast a perturbation decays are different questions with different answers.

\subsection*{Supplementary Methods}\label{sec:smethods}

\subsubsection*{Conventions}

We write the AR($p$) model as in eq. \ref{eq:arp}, so that the characteristic polynomial is $z^p - \varphi_1 z^{p-1} - \cdots - \varphi_p$ and stationarity requires all of its roots to lie inside the unit circle. Sign conventions differ between texts, and the two common polynomials are reciprocals of one another; we use this one throughout, so that a pole inside the unit circle means a stationary process. The one-sided spectral density is
\begin{align}
S(f) = \frac{2\sigma^2/f_s}{\left|1 - \sum_{k=1}^{p}\varphi_k e^{-i 2\pi f k/f_s}\right|^{2}}, \qquad 0 \le f \le f_s/2,
\end{align}
normalized so that its integral is the process variance, matching Welch's method. For $p=1$ this is the discrete-time Lorentzian, whose knee frequency is $1/(2\pi\tau)$. The theoretical autocovariance is obtained by solving the Yule-Walker system for the first $p+1$ lags and recursing thereafter. A fit inherited from elsewhere with a pole at $|\lambda| \ge 1$ can be made stationary by reflecting the pole to $1/\bar\lambda$ and dividing the scale by $|\lambda|^2$, which leaves the model spectrum unchanged.

\subsection*{Supplementary Figures}

\begin{figure}[htbp]
  \centering
  \includesvg[width=\linewidth]{revision/figures/figS_oscillation_grid.svg}
  \caption{An added sinusoid, swept over frequency and power. Median estimated aperiodic timescale divided by the truth, across oscillation frequency and oscillation power, for the four fitted estimators. One is exact; the colour scale is logarithmic, so a factor of two too large and a factor of two too small are equally far from white. Every cell is filled because the fit is constrained to contain a relaxation mode; before that constraint the lag-domain fits left cells empty, having spent every pole on the rhythm. Both spectral fits stay between 0.97 and 1.03 everywhere: robust regression discards the rhythm and leaves the aperiodic fit untouched.}
  \label{fig:s_grid}
\end{figure}

\begin{figure}[htbp]
  \centering
  \includesvg[width=\linewidth]{revision/figures/figS_windows.svg}
  \caption{Real recordings cut into windows. \textbf{a) Convergence.} Each recording cut into non-overlapping windows, the estimate normalized to that from the longest window and plotted against window length in units of that estimate; shading is $\pm 20\%$. \textbf{b) Excess spread.} Spread of estimates across windows of a given length, divided by the spread finite sampling alone predicts, at the longest window available. Values above one are variability that sampling cannot explain.}
  \label{fig:s_windows}
\end{figure}

\clearpage

\bibliography{references,references_revision}


\end{document}
